\chapter{Về Việc Học}
\markboth{Về Việc Học}{On Learning}

\section{Học không phải để hơn người khác, mà để hơn chính mình hôm qua.}
\begin{truyen}{Hơn chính mình hôm qua}{Better Than Yesterday's Self}
\chuhoa{M}{ột học sinh hỏi thầy: ``Thưa thầy, làm sao em biết em đã giỏi?'' Thầy không so em với ai trong lớp. Thầy chỉ hỏi: ``Hôm qua em chưa làm được bài này, hôm nay em làm được. Đó chính là giỏi.''}
\textit{(A student asked: ``How do I know I'm good?'' The teacher didn't compare him to anyone. He only said: ``Yesterday you couldn't do this; today you can. That is being good.'')}

Từ đó em ấy bớt nhìn sang bàn bên, và bắt đầu nhìn lại vở mình — xem tuần trước mình viết thế nào, hôm nay đã rõ hơn chưa.
\textit{(From then on, that student looked less at the next desk and more at his own notebook — to see how he wrote last week and whether today was a bit clearer.)}
\end{truyen}

\begin{baihoc}
Thước đo quan trọng nhất không phải bạn bè, mà là chính mình của ngày hôm qua.
\textit{(The most important measure is not your peers, but yesterday's you.)}
\end{baihoc}

\begin{ghinhoanh}
\textbf{Short English to remember:}

Learn to be better than who you were yesterday, not better than someone else.

\end{ghinhoanh}

\begin{tuhocnhanh}
1. Vì sao so với chính mình hôm qua lại có ích? \textit{Why is comparing with yesterday's self useful?}

2. Em có thể tự đánh giá tiến bộ của mình bằng cách nào? \textit{How can you measure your own progress?}
\end{tuhocnhanh}
\ngancach

\section{Điểm số nói lên kết quả, nhưng thái độ nói lên tương lai.}
\begin{truyen}{Điểm và thái độ}{Grades and Attitude}
\chuhoa{C}{ó hai học sinh cùng điểm trung bình. Một em mỗi lần trả bài đều gật đầu lắng nghe, ghi lại chỗ sai để xem lại. Em kia vừa nhận bài đã gấp vội bỏ cặp. Thầy nói với cả lớp: ``Điểm bây giờ có thể giống nhau, nhưng vài năm nữa hai con đường sẽ khác nhau — vì thái độ khác nhau.''}
\textit{(Two students had the same average. One always listened carefully when papers were returned and wrote down mistakes to review. The other folded the paper and put it away. The teacher told the class: ``Grades may be the same now, but in a few years the two paths will differ — because attitudes differ.'')}
\end{truyen}

\begin{baihoc}
Thái độ với bài sai và với việc học quyết định em đi xa đến đâu, không chỉ điểm một kỳ.
\textit{(Your attitude toward mistakes and learning determines how far you go, not just one semester's grade.)}
\end{baihoc}

\begin{ghinhoanh}
\textbf{Short English to remember:}

Grades show results; attitude shows the future.

\end{ghinhoanh}

\begin{tuhocnhanh}
1. Thái độ nào với điểm và với bài sai là tốt? \textit{What attitude toward grades and mistakes is good?}

2. Em có đang chú ý rèn thái độ hay chỉ chạy theo điểm? \textit{Do you work on attitude or only on grades?}
\end{tuhocnhanh}
\ngancach

\section{Một ngày không học được gì mới, là một ngày mất đi một cơ hội trưởng thành.}
\begin{truyen}{Một ngày mới}{A New Day}
\chuhoa{T}{hầy hay hỏi cuối giờ: ``Hôm nay em học được điều gì mới — dù chỉ một ý nhỏ?'' Nếu em nào nói ``dạ không ạ'', thầy không la. Thầy chỉ nói nhẹ: ``Vậy thì ngày mai mình cố tìm lấy một ý. Mỗi ngày một ý nhỏ, cả năm đã là cả trăm ý.''}
\textit{(The teacher often asked at the end of class: ``What new thing did you learn today — even one small idea?'' If someone said ``nothing'', he didn't scold. He just said softly: ``Then tomorrow let's try to find one. One small idea a day, and in a year that's hundreds.'')}
\end{truyen}

\begin{baihoc}
Trưởng thành không cần nhảy vọt; mỗi ngày một chút mới thật sự tích lũy.
\textit{(Growth doesn't need a leap; a little that is new each day truly adds up.)}
\end{baihoc}

\begin{ghinhoanh}
\textbf{Short English to remember:}

A day without learning something new is a day without a chance to grow.

\end{ghinhoanh}

\begin{tuhocnhanh}
1. Hôm nay em học được điều gì mới? \textit{What did you learn new today?}

2. Làm sao để mỗi ngày đều có ít nhất một ý mới? \textit{How can you have at least one new idea every day?}
\end{tuhocnhanh}
\ngancach

\section{Người học giỏi không phải người biết nhiều nhất, mà là người hỏi nhiều nhất.}
\begin{truyen}{Người hỏi nhiều}{The One Who Asks the Most}
\chuhoa{L}{ớp có em hay giơ tay hỏi, có khi hỏi đến mức bạn bảo ``sao hỏi nhiều thế''. Thầy nói trước lớp: ``Người giỏi không phải người trả lời đúng nhiều nhất, mà người dám hỏi nhiều nhất. Vì hỏi là đang nghĩ; không hỏi đôi khi là không dám nghĩ.''}
\textit{(One student often raised his hand to ask, so much that friends said ``why so many questions''. The teacher told the class: ``The strong one is not who answers correctly the most, but who dares to ask the most. Because to ask is to think; not asking sometimes means not daring to think.'')}
\end{truyen}

\begin{baihoc}
Câu hỏi là dấu hiệu của tư duy; im lặng mãi có thể là im vì sợ sai.
\textit{(Questions are a sign of thinking; long silence can mean silence out of fear of being wrong.)}
\end{baihoc}

\begin{ghinhoanh}
\textbf{Short English to remember:}

The best learners ask the most questions.

\end{ghinhoanh}

\begin{tuhocnhanh}
1. Vì sao hỏi nhiều lại là dấu hiệu học giỏi? \textit{Why is asking a lot a sign of learning well?}

2. Em có dám hỏi khi chưa hiểu không? \textit{Do you dare to ask when you don't understand?}
\end{tuhocnhanh}
\ngancach

\section{Sai không đáng sợ, sợ nhất là không dám thử.}
\begin{truyen}{Dám thử}{Daring to Try}
\chuhoa{M}{ột em không bao giờ xung phong lên bảng vì sợ sai. Thầy gọi em lên và nói: ``Thầy không trừ điểm khi em làm sai. Thầy chỉ buồn khi em không thử. Lần này em thử đi — sai thì mình sửa, cả lớp cùng học.''}
\textit{(A student never volunteered to go to the board for fear of being wrong. The teacher called him up and said: ``I don't take points off for wrong answers. I'm only sad when you don't try. This time just try — if it's wrong we fix it, the whole class learns together.'')}
\end{truyen}

\begin{baihoc}
Không thử thì không bao giờ biết mình sai ở đâu để sửa; không sửa thì không tiến.
\textit{(If you don't try, you never know where you're wrong to fix it; if you don't fix, you don't move forward.)}
\end{baihoc}

\begin{ghinhoanh}
\textbf{Short English to remember:}

Being wrong is not the worst; not daring to try is.

\end{ghinhoanh}

\begin{tuhocnhanh}
1. Điều gì đáng sợ hơn việc làm sai? \textit{What is scarier than making a mistake?}

2. Em có thường né thử vì sợ sai không? \textit{Do you often avoid trying because you're afraid of being wrong?}
\end{tuhocnhanh}
\ngancach

\section{Nếu hôm nay em hiểu hơn hôm qua một chút, em đã tiến bộ rồi.}
\begin{truyen}{Một chút hiểu thêm}{A Little More Understanding}
\chuhoa{E}{m hay so mình với bạn giỏi nhất lớp nên lúc nào cũng thấy mình kém. Thầy gọi em ra nói: ``Em đừng so với người đứng đầu. Em so với chính em: tuần trước em chưa hiểu đoạn này, hôm nay em đã hiểu. Đó là tiến bộ thật.''}
\textit{(A student always compared himself to the top of the class and felt behind. The teacher called him aside: ``Don't compare yourself to the first. Compare with yourself: last week you didn't get this part, today you do. That's real progress.'')}
\end{truyen}

\begin{baihoc}
Tiến bộ là so với chính mình: hôm qua chưa hiểu, hôm nay hiểu — dù chỉ một chút.
\textit{(Progress is measured against yourself: yesterday you didn't understand, today you do — even just a little.)}
\end{baihoc}

\begin{ghinhoanh}
\textbf{Short English to remember:}

If today you understand a little more than yesterday, you have already improved.

\end{ghinhoanh}

\begin{tuhocnhanh}
1. Em có hay so mình với người giỏi nhất không? \textit{Do you often compare yourself to the best?}

2. Làm sao để thấy rõ tiến bộ của chính mình? \textit{How can you see your own progress clearly?}
\end{tuhocnhanh}
\ngancach

\section{Kiến thức giống như tiền tiết kiệm: mỗi ngày góp một ít, sau này sẽ rất giàu.}
\begin{truyen}{Góp mỗi ngày}{Saving a Little Each Day}
\chuhoa{T}{hầy ví việc học như bỏ ống heo: ``Mỗi ngày em học thuộc một công thức, hiểu thêm một định nghĩa — giống như mỗi ngày bỏ vào ống một đồng. Không thấy ngay, nhưng vài năm sau mở ra, em sẽ thấy mình đã giàu kiến thức.''}
\textit{(The teacher compared learning to a piggy bank: ``Every day you memorize one formula, understand one more definition — like putting one coin in each day. You don't see it at once, but open it in a few years and you'll see you're rich in knowledge.'')}
\end{truyen}

\begin{baihoc}
Học đều mỗi ngày một chút bền hơn học dồn một lúc; kiến thức cũng cần tích lũy.
\textit{(Learning a little every day lasts longer than cramming once; knowledge needs to accumulate too.)}
\end{baihoc}

\begin{ghinhoanh}
\textbf{Short English to remember:}

Knowledge is like savings: a little each day makes you rich later.

\end{ghinhoanh}

\begin{tuhocnhanh}
1. Vì sao ``góp mỗi ngày'' tốt hơn ``học dồn''? \textit{Why is ``a little each day'' better than cramming?}

2. Em đang góp đều hay chỉ góp khi gần thi? \textit{Are you saving steadily or only before exams?}
\end{tuhocnhanh}
\ngancach

\section{Học giỏi không phải vì thông minh, mà vì không bỏ cuộc.}
\begin{truyen}{Không bỏ cuộc}{Not Giving Up}
\chuhoa{C}{ó em nghĩ mình không thông minh nên học mãi không lên. Thầy kể: ``Nhiều người giỏi thật ra không phải IQ cao nhất. Họ chỉ là người làm lại bài sai nhiều lần, hỏi đến khi hiểu, không bỏ cuộc khi khó. Kiên trì mới là thứ làm nên sự khác biệt.''}
\textit{(A student thought he wasn't smart so his grades wouldn't improve. The teacher said: ``Many top students don't have the highest IQ. They're just the ones who redo wrong problems, ask until they understand, and don't quit when it's hard. Persistence is what makes the difference.'')}
\end{truyen}

\begin{baihoc}
Thông minh giúp đi nhanh; không bỏ cuộc giúp đi xa.
\textit{(Smarts help you go fast; not giving up helps you go far.)}
\end{baihoc}

\begin{ghinhoanh}
\textbf{Short English to remember:}

Doing well comes from not giving up, not from being the smartest.

\end{ghinhoanh}

\begin{tuhocnhanh}
1. Điều gì quan trọng hơn thông minh trong việc học? \textit{What matters more than smarts in learning?}

2. Em có hay bỏ cuộc khi gặp bài khó không? \textit{Do you often give up when a problem is hard?}
\end{tuhocnhanh}
\ngancach

\section{Người lười học hôm nay sẽ phải học rất nhiều từ cuộc đời sau này.}
\begin{truyen}{Học từ cuộc đời}{Learning from Life Later}
\chuhoa{T}{hầy nói với một em hay trốn học, chơi game: ``Bây giờ em không chịu học bài trong sách, sau này cuộc đời sẽ dạy em bằng cách khác — bằng thất bại, bằng thiếu cơ hội, bằng hối tiếc. Bài học lúc ấy đắt hơn và đau hơn nhiều.''}
\textit{(The teacher told a student who often skipped class to play games: ``If you won't learn from books now, life will teach you later in another way — through failure, missed chances, regret. Those lessons cost more and hurt more.'')}
\end{truyen}

\begin{baihoc}
Học trong lớp là con đường rẻ nhất và an toàn nhất; trốn học hôm nay là trả giá đắt sau này.
\textit{(Learning in class is the cheapest and safest path; skipping it today means paying a higher price later.)}
\end{baihoc}

\begin{ghinhoanh}
\textbf{Short English to remember:}

Those who don't learn today will have to learn much more from life later.

\end{ghinhoanh}

\begin{tuhocnhanh}
1. Vì sao học ở trường lại ``rẻ'' hơn học từ cuộc đời? \textit{Why is learning at school ``cheaper'' than learning from life?}

2. Em có đang trốn học bằng cách nào không? \textit{Are you avoiding learning in any way?}
\end{tuhocnhanh}
\ngancach

\section{Điều quan trọng không phải là em học bao lâu, mà là em học có suy nghĩ hay không.}
\begin{truyen}{Học có suy nghĩ}{Learning with Thinking}
\chuhoa{M}{ột em ngồi học ba tiếng nhưng toàn nhìn sách, không ghi chú, không tự hỏi. Thầy nói: ``Ba tiếng ngồi im mà không suy nghĩ, không bằng ba mươi phút mà em đọc — rồi đặt câu hỏi, tóm tắt, tự giải thích lại. Thời gian không quyết định; cách học mới quyết định.''}
\textit{(A student sat for three hours but only stared at the book — no notes, no questions. The teacher said: ``Three hours of sitting without thinking is worth less than thirty minutes where you read, then ask questions, summarize, explain back. Time doesn't decide; how you learn does.'')}
\end{truyen}

\begin{baihoc}
Học có suy nghĩ là học có hỏi, có tóm tắt, có liên hệ — không phải chỉ mắt đọc qua trang.
\textit{(Learning with thinking means asking, summarizing, connecting — not just eyes passing over the page.)}
\end{baihoc}

\begin{ghinhoanh}
\textbf{Short English to remember:}

What matters is not how long you study, but whether you think while you learn.

\end{ghinhoanh}

\begin{tuhocnhanh}
1. Học ``có suy nghĩ'' nghĩa là làm những gì? \textit{What does ``learning with thinking'' mean in practice?}

2. Em thường học theo kiểu thời gian hay kiểu suy nghĩ? \textit{Do you usually learn by clock time or by thinking?}
\end{tuhocnhanh}
\ngancach
